% Working paper-selection document. Replace with a four-page ISBI report
% after choosing one paper and completing its experiments.
\documentclass{article}
\usepackage{ISBI_template-master/spconf,amsmath,amssymb,graphicx}
\usepackage{booktabs,longtable,array,enumitem}
\newcolumntype{P}[1]{>{\raggedright\arraybackslash}p{#1}}
\setlength{\parindent}{0pt}
\setlist[itemize]{nosep,leftmargin=1.4em}
\pagestyle{plain}

\title{Assignment 2: Feasible Paper Reproduction Plans}
\name{Group paper-selection working document}
\address{Advanced Topics in Deep Learning, University of Copenhagen}

\begin{document}
\onecolumn
\maketitle

\section{Assignment and the two tracks}
The group selects \emph{one} curated paper from Topics 1--4. Each group
member must reproduce a distinct experiment. Only Topics 1--3, comprising
18 curated papers, have been published in the local course materials so far.
This comparison is for paper selection, not the final four-page submission.

\textbf{Baselines track:} reproduce and evaluate comparison methods
\emph{actually evaluated in the selected paper}, with relevant ablations or
hyperparameter tuning. The paper's new method need not be retrained. A
control of that method, or a different model merely cited in related work,
does not by itself make a baselines-track experiment. Papers without an
attainable reported comparator are marked ``no strict plan'' below; the
listed control studies for those papers require the instructor to accept a
broader interpretation of the track.

\textbf{Replication track:} reproduce a named subset of the paper's main
results using its released code when available, or a faithful
reimplementation, and include a targeted ablation in the costed core.
Released-weight evaluation is allowed by the assignment and can replicate a
reported evaluation table, but it cannot establish a training-recipe claim.
Such plans are labeled \emph{evaluation-only} below. Smaller models, datasets,
resolutions, or sample counts must be disclosed, together with their likely
effect on the conclusion.

\section{Cost model and scope}
Every total below is an \emph{alternative plan for one paper and one track},
not a total across all 18 papers. Where a track fits, its core plan gives
three distinct contributions for a three-person group; a fourth member can
take an optional evaluation from the same paper. The menus collect feasible
experiment families, not an instruction to run every combination. The
\$1,000 ceiling applies to the selected plan. No large generative or
foundation-model pretraining is budgeted; AnyUp's reported five-hour
feature-upsampler training is an explicit small-scale exception.

GPU-hours count devices: two GPUs for six elapsed hours are 12 GPU-h. For each
component we estimate runs $\times$ B200 GPUs $\times$ elapsed hours, then
multiply by Prime Intellect's publicly listed \textbf{B200 rate of
\$3.49 per GPU-hour}, checked 23 September 2026. The listing shows a
single-GPU option; market prices and availability can change. The run-hour
allowances are planning placeholders, largely carried over from earlier
mixed-GPU estimates and adjusted for the track corrections below. No
unmeasured B200 speedup is credited.
They are \emph{not} benchmarked B200 runtimes.
Shared feature extraction, generation, and scoring appear once per track.
The totals exclude CPU, RAM, storage, data transfer, human annotation, paid
model APIs, and engineering time. CPU-heavy cases are marked explicitly.
Treat every runtime as an \emph{untimed assumption}, except where the paper
reports a specific throughput; run a 20--100-step or image pilot on the
intended remote GPU before committing. Reserve up to \textbf{2$\times$ the
quoted GPU amount} for slower throughput, reruns, and failed setup. Even the
largest proposed core plan, REPA replication at \$383.90, remains below
\$1,000 at 2$\times$ (\$767.80). These are GPU rental estimates, not guarantees of an
exact reproduction.

\clearpage
\section{Overview: one three-person core plan per track}
Costs in this table are rounded to the nearest dollar. ``N/A'' means there is
no strict baselines-track plan at a defensible GPU-only cost; conditional
control studies below are not included as compliant plans. The itemized
calculations below are the actual comparison basis.
\begingroup\footnotesize
\setlength{\tabcolsep}{2.5pt}\renewcommand{\arraystretch}{1.18}
\begin{longtable}{@{}P{0.34in}P{0.65in}P{1.48in}P{1.34in}P{0.47in}P{1.34in}P{0.47in}@{}}
\toprule
Week & Topic & Paper & Baselines core & GPU cost & Replication core & GPU cost \\
\midrule
\endfirsthead
\toprule
Week & Topic & Paper & Baselines core & GPU cost & Replication core & GPU cost \\
\midrule
\endhead
1--2 & Theory & Opening the Black Box of Deep Neural Networks via Information & No distinct reported comparator & N/A & Information paths and binning ablation & \$26 \\
1--2 & Theory & Computing Nonvacuous Generalization Bounds & SGD error, VC and Rademacher bounds & \$15 & Posterior bound and variance ablation & \$108 \\
1--2 & Theory & A Bayesian Perspective on Generalization and SGD & No distinct reported comparator & N/A & Batch, learning-rate, dataset-size sweeps & \$42 \\
1--2 & Theory & Deep Double Descent & No distinct reported comparator & N/A & Width, sample, and label-noise curves & \$195 \\
1--2 & Theory & Towards Understanding Grokking & No distinct reported comparator & N/A & Addition, modular transformer, ablation & \$48 \\
1--2 & Theory & Binarized Neural Networks Converge Toward Algorithmic Simplicity & Entropy versus BDM & \$21 & Width, depth, structured-data controls & \$38 \\
2--3 & Generative & Representation Alignment for Generation (REPA) & Unaligned DiT and SiT & \$244 & SiT control, REPA, alignment ablation & \$384 \\
2--3 & Generative & Mean Flows for One-step Generative Modeling & Flow baselines and NFE curve & \$216 & MeanFlow, control, loss ablation & \$307 \\
2--3 & Generative & Diffusion Models for Counterfactual Explanations (DiME) & DiVE and DiVE$^{100}$ comparisons & \$154 & DiME, DiVE, guidance ablation & \$154 \\
2--3 & Generative & DiffAtlas & Segmentation and diffusion baselines & \$140 & Full, few-shot, guidance ablation & \$181 \\
2--3 & Generative & FLUX.1 Kontext & No GPU-only reported comparator plan & N/A & [dev] benchmark and chain ablation; evaluation-only & \$91 \\
2--3 & Generative & Qwen-Image & Generator and VAE baselines & \$140 & Benchmark subsets and guidance ablation; evaluation-only & \$154 \\
4--5 & Foundation & DINOv3 & DINOv2/SigLIP2 frozen probes & \$70 & Transfer probes; evaluation-only & \$63 \\
4--5 & Foundation & V-JEPA 2.1 & V-JEPA 2/DINOv2 probes & \$112 & Video and dense probes; evaluation-only & \$112 \\
4--5 & Foundation & SAM 3 & Concept and interactive baselines & \$63 & Concept prompts and video; evaluation-only & \$77 \\
4--5 & Foundation & AnyUp & Bilinear, FeatUp, LoftUp, JAFAR & \$84 & Train and ablate AnyUp & \$105 \\
4--5 & Foundation & Med-SA & SAM/nnU-Net/MedSAM baselines & \$154 & Adapter, slice, prompt ablations & \$307 \\
4--5 & Foundation & Unified View of Parameter-Efficient Transfer Learning & PEFT method comparisons & \$181 & Adapter placement, scaling, MAM & \$168 \\
\bottomrule
\end{longtable}
\endgroup

\section{Topic 1: Deep Learning Theory, weeks 1--2}

\subsection{Opening the Black Box of Deep Neural Networks via Information}
The paper uses all 4,096 possible 12-bit inputs, a near-balanced stochastic
spherical-harmonic label rule, a narrow tanh MLP, SGD, 30 activation bins, and
50 original random runs. A parity or XOR substitute would be a different task.
The paper does not evaluate an independent comparator method. Its depth,
sample-fraction, and teacher-rule conditions are method controls, so they do
not form a strict baselines-track plan.

\textbf{Conditional control study, not a baselines-track core:} (1) Depth versus epochs to generalization, four
depths $\times$ two seeds $\times$ 0.3 B200-h = 2.4 h; (2) 5\%, 45\%, and 85\%
training fractions, three $\times$ two $\times$ 0.3 h = 1.8 h; (3) symmetric
versus nonsymmetric teacher, two additional 0.3 h runs = 0.6 h. Total \textbf{4.8 B200-h $\times$ \$3.49 = \$16.75}. Other feasible checks: 15/30/60-bin
mutual-information estimates from saved activations (CPU), and test-error
curves for each depth.

\textbf{Replication core:} (1) Information-plane fit-then-compress trajectories
at early/middle/late checkpoints, three seeds $\times$ 0.75 h = 2.25 B200-h;
(2) depth effect, four networks $\times$ 0.75 h = 3 h; (3) sample-fraction
effect, three $\times$ 0.75 h = 2.25 h. The targeted ablation recomputes the
information trajectories with 15/30/60 activation bins from saved activations
on CPU. Total \textbf{7.5 B200-h $\times$ \$3.49 = \$26.18}, plus CPU. From the
same checkpoints, gradient mean/SD and signal-to-noise transition, an IB
reference curve, and layerwise fixed-point tests are also feasible. The
reported ``compression'' is sensitive to discretization; retain the paper's
30-bin protocol for the primary comparison.

\subsection{Computing Nonvacuous Generalization Bounds for Deep (Stochastic) Neural Networks with Many More Parameters than Training Data}
Use binary MNIST digits 0--4 versus 5--9, the original 55k training split,
and a data-independent initialization prior. The true-label bound uses a
diagonal Gaussian posterior around an SGD solution; the original optimizes it
for 200k steps and uses a large Monte Carlo evaluation. This paper is
methodologically risky even though its models are small: a bound is only
meaningful if its prior, data split, KL inversion, and sampling correction are
valid.

\textbf{Baselines core:} (1) Deterministic SGD test error for three reported
MLP widths/depths, three $\times$ 0.5 B200-h = 1.5 h; (2) reproduce the
paper's VC and path-norm/Rademacher comparison bounds on those networks, with
three regularization settings, three $\times$ 0.3 h = 0.9 h plus CPU bound
calculation; (3) random-label memorization as the reported vacuity control,
2 h. Total \textbf{4.4 B200-h $\times$ \$3.49 = \$15.36}, plus CPU. Compare each
bound with the same network's test error; deterministic-versus-stochastic
test error is another feasible check.

\textbf{Replication core:} (1) 600-unit true-label posterior and certified
bound, 8 B200-h; (2) random-label posterior as a vacuity control, 15 h;
(3) a posterior-variance ablation at the same
architecture, 8 h. Total \textbf{31 B200-h $\times$ \$3.49 = \$108.19}. A
wider/deeper posterior is an optional 10 h or \$34.90. A limited run may use
minibatches and 1k--5k independent posterior
draws, with the appropriate confidence correction, but nonvacuity is not
guaranteed. Do not tune an allegedly data-independent prior on labels without
the paper's correction.

\subsection{A Bayesian Perspective on Generalization and SGD}
The paper has two attainable settings: an 800-example MNIST 0-versus-1
logistic model for approximate evidence, and an 800-hidden-unit MNIST MLP on
1,000 examples trained for 10k updates. Compare at fixed update counts rather
than silently giving large batches more data passes. These are the paper's
own analyses rather than comparisons with a distinct reported baseline.

\textbf{Conditional control study, not a baselines-track core:} (1) True-versus-random-label logistic regression and
regularization/evidence sweep (CPU); (2) MLP test accuracy across six batch
sizes, six $\times$ 0.4 B200-h = 2.4 h; (3) L2 regularization at two batch
sizes, two $\times$ 0.4 h = 0.8 h. Total \textbf{3.2 B200-h $\times$ \$3.49 = \$11.17}, plus CPU. Minibatch-gradient histograms and train/test
cross-entropy are further cheap checks.

\textbf{Replication core:} (1) Six-batch optimum curve, 2.4 B200-h;
(2) three-learning-rate by four-batch grid, 4.8 h; (3) three-training-size by
four-batch grid, 4.8 h. The learning-rate and sample-size changes are targeted
ablations of the predicted effective noise scale. Total \textbf{12 B200-h $\times$ \$3.49 = \$41.88}.
Momentum-by-batch is another 4.8 h (\$16.75); logistic Laplace-evidence versus
held-out loss is CPU work. The claimed optimum shifts with effective SGD
noise scale, approximately $\epsilon N/[B(1-m)]$; preserve the paper's
learning-rate-dependent step normalization when testing that claim.

\subsection{Deep Double Descent}
The paper varies CNN/ResNet widths, training epochs, sample count, label
noise, augmentation, and optimizer. The feasible core uses a 5k--10k CIFAR-10
subset and five-layer CNN widths chosen around the \emph{observed}
interpolation threshold. Budget about four B200-h per completed width run, then
adjust after a pilot. If near-threshold models never reach near-zero training
error, the central double-descent claim has not been tested.
Random features, label noise, and optimizer changes are reported conditions
of the phenomenon, not baseline methods for a new technique.

\textbf{Conditional control study, not a baselines-track core:} (1) Random-feature width/sample sweep on
Fashion-MNIST (CPU); (2) five CNN widths at zero and 15\% fixed label noise,
ten $\times$ 4 = 40 B200-h; (3) SGD versus Adam at three widths, six $\times$
4 = 24 h. Total \textbf{64 B200-h $\times$ \$3.49 = \$223.36}; shared control
runs can lower the bill. Other feasible contrasts: augmentation on/off at
three widths (24 h, \$83.76), weight decay/early stopping (about 12 h,
\$41.88), five-model ensembling near interpolation (20 h, \$69.80), and a
small-transformer translation sweep (roughly 120 B200-h, \$418.80; high
setup burden). The translation sweep is an alternative project, not an
addition to every core run under the 2$\times$ reserve.

\textbf{Replication core:} (1) Seven-width model-wise double-descent curve,
seven $\times$ 4 = 28 B200-h, with epoch-wise curves from saved checkpoints;
(2) four sample sizes at one critical width, four $\times$ 4 = 16 h;
(3) a targeted label-noise ablation near that width, three $\times$ 4 = 12 h.
Total \textbf{56 B200-h $\times$ \$3.49 = \$195.44}. Optional tests include
augmentation or optimizer shifts (about 12 h/\$41.88 each) and a four-width adversarial-robustness
sweep (about 60 B200-h/\$209.40). These are reduced-scale trends, not the
original 4,000-epoch or 500k-step grids.

\subsection{Towards Understanding Grokking}
Keep separate the paper's small commutative addition task, permutation-group
$S_3$ task, and $p=53$ modular-addition transformer. The first uses unordered
held-out pairs. This theory paper has reported control conditions but no
distinct evaluated comparison method.

\textbf{Conditional control study, not a baselines-track core:} (1) Toy addition decoder/embedding learning-rate and
weight-decay grid, nine tiny runs $\times$ 0.1 B200-h = 0.9 h; (2) $S_3$ phase
comparison, six $\times$ 0.1 = 0.6 h; (3) representation-quality index versus
held-out accuracy at three data fractions and three seeds, nine $\times$
0.1 = 0.9 h. Total \textbf{2.4 B200-h $\times$ \$3.49 = \$8.38}; CPU may be
preferable. Dropout on/off for the modular transformer is another 3 B200-h or
\$10.47.

\textbf{Replication core:} (1) Toy addition memorization-to-grokking curves
at three fractions and three seeds, 1.8 B200-h; (2) $p=53$ modular-transformer
curves and embedding geometry, three $\times$ 2 = 6 h; (3) weight-decay or
dropout ablation, three $\times$ 2 = 6 h. Total \textbf{13.8 B200-h $\times$ \$3.49 = \$48.16}. Predicted accuracy from parallelogram structure,
representation-quality index, and effective-theory grokking time can be
computed from saved checkpoints on CPU. A small MNIST phase plot is another
possible, less directly specified evaluation.

\subsection{Binarized Neural Networks Converge Toward Algorithmic Simplicity}
The paper uses small STE-binarized MLPs, 10-by-10 resized MNIST, BDM from
4-by-4 sign-weight blocks, matched Shannon entropy, 200 independent subsets
per original condition, and bootstrapped loss-complexity correlations. We
propose 20 seeds per condition and must report wider uncertainty. BDM lookup
is CPU-heavy; the following B200 rental is a billed-instance proxy, not a
claim that GPU arithmetic dominates.

\textbf{Baselines core:} reproduce the paper's Shannon-entropy comparator
against block-decomposition-method (BDM) complexity using the same saved
weights and matched partitions: (1) three MNIST widths, 20 seeds each, about
3 B200-h; (2) shuffled-label and random-input controls, 20 each, 2 h;
(3) Fashion-MNIST transfer, 20 seeds, 1 h. Total \textbf{6 B200-h $\times$ \$3.49 = \$20.94}, plus CPU for entropy, BDM, correlations, and
bootstrap intervals. The width and dataset changes test whether the reported
entropy-versus-BDM ranking survives these settings; checkpoint cadence is a
further cheap sensitivity check.

\textbf{Replication core:} (1) MNIST BDM-minus-entropy width trend, 3 B200-h;
(2) one- to four-layer depth trend at matched parameter count, 4 h;
(3) structured MNIST/Fashion/synthetic-signal versus shuffled/random controls,
4 h. Total \textbf{11 B200-h $\times$ \$3.49 = \$38.39}, plus CPU. UCI HAR
engineered-versus-raw features add about 2 h (\$6.98). Use the same weight
partition for BDM and entropy. Pilot CPU throughput before renting any GPU;
the original largest 200-run configurations took about 15 hours on an M4 Max.

\section{Topic 2: Generative Models, weeks 2--3}

\subsection{Representation Alignment for Generation: Training Diffusion Transformers Is Easier Than You Think (REPA)}
The central question is whether alignment with a frozen image encoder helps a
diffusion transformer learn faster. Use one public image dataset/subset at
about 128 pixels, fixed VAE latents and DINOv2 features, matched small SiT or
DiT backbones, and identical optimization steps. Feasible evaluations include
quality versus checkpoint, linear probes or representation similarity,
2k--5k-sample relative FID, sample precision/recall, alignment-layer and
loss-weight ablations, encoder choice (DINOv2/CLIP/MAE), and sampler
steps/classifier-free guidance. With fewer than the paper's 50k samples and
250 denoising calls per sample, FID is a noisy \emph{relative} measure, not
comparable numerically with its headline result.

\textbf{Baselines core:} (1) Small DiT without alignment, 30 B200-h;
(2) matched SiT without alignment, 30 h; (3) checkpoint-quality and
sampler/guidance comparison on the same fixed samples, 10 h. Total \textbf{70 B200-h $\times$ \$3.49 = \$244.30}. An optional sampler setting
on the unaligned comparators uses shared checkpoints and adds roughly
5--15 B200-h. Changing REPA's alignment encoder belongs to replication.

\textbf{Replication core:} (1) SiT no-REPA control, 30 B200-h; (2) matched
SiT+REPA, 34 h; (3) one changed alignment depth or target encoder, 34 h;
shared feature caching, 2k--5k image generation at a shorter fixed sampler,
and representation probes, 12 h. Total \textbf{110 B200-h $\times$ \$3.49 = \$383.90}. A loss-weight or NT-Xent-versus-cosine ablation is another
roughly 34 B200-h/\$118.66. The paper reports eight H100s at 5.4 steps/s:
one 400k-step run is about $(400000/5.4/3600)\times8=165$ \emph{H100}
GPU-h. If a B200 run took the same GPU-hours, that would be about
165 $\times$ \$3.49 = \$575.85 before evaluation; a matched pair would exceed
\$1,000. Actual B200 throughput must be measured. The proposed plan tests
a reduced-scale trend only.

\subsection{Mean Flows for One-step Generative Modeling}
Use a compact CIFAR-10 32-pixel U-Net or latent ViT and fixed 50k--100k
training steps; compare at identical data and model size. Feasible questions
include flow matching ($r=t$) versus MeanFlow, one/two/four/eight function
evaluations (NFE), generation latency, nontrivial $r\ne t$ sampling ratio,
loss exponent, time-pair distribution, conditioning representation, guidance
strength, and small-width scaling. Use about 10k generated samples for a
relative FID; it is not the paper's FID-50K.

\textbf{Baselines core:} (1) Matched small flow-matching training, 22 B200-h;
(2) second time-sampling or hyperparameter baseline, 22 h; (3) evaluate a
released iCT/ECT/Shortcut baseline if obtainable, 10 h; shared
one/two/four-NFE sampling and FID, 8 h. Total \textbf{62 B200-h $\times$ \$3.49 = \$216.38}. If those checkpoints are unavailable, use a more complete
flow-matching NFE curve from the same control instead and report that the
published comparator could not be reproduced.

\textbf{Replication core:} (1) Flow-matching control, 22 B200-h;
(2) MeanFlow default, 28 h; (3) loss exponent or $r\ne t$ ratio ablation,
28 h; shared sampling/FID/latency, 10 h. Total \textbf{88 B200-h $\times$ \$3.49 = \$307.12}. An incorrect-JVP-target or alternate time-conditioning
ablation adds about 28 B200-h/\$97.72. The paper's B/4 ablation uses 400k ImageNet
steps, and its timing is on TPU v4-8; neither is a measured GPU runtime for
this smaller scenario.

\subsection{Diffusion Models for Counterfactual Explanations (DiME)}
Use CelebA-128, a fixed attribute classifier and diffusion checkpoint, and
the \emph{same} 100--200 inputs per attribute across methods. Feasible
evaluations are smile versus young flips; validity/flip rate, identity (FVA),
minimality (L1/LPIPS/MNAC), realism, five-seed diversity, spurious-correlation
difference, and guidance-strength/noise-level/step-count ablations. Original
sampling uses about 1,800 denoiser calls per explanation: 100 examples imply
180k calls. At an illustrative 0.05 s/call that is 2.5 GPU-h before retries,
gradients, and metrics, motivating a roughly 10 h slot per condition.

\textbf{Baselines core:} reproduce the paper's Table 1 comparator DiVE rather
than its Table 2 DiME variations: (1) DiVE on smile, 8 B200-h; (2) DiVE$^{100}$
with 100 optimization steps on smile, 12 h; (3) both settings on young,
8+12 = 20 h;
shared scoring/oracles and a small optimization-step tuning check, 4 h. Total
\textbf{44 B200-h $\times$ \$3.49 = \$153.56}. Report flip rate, identity FVA,
MNAC, and diversity on the same inputs. xGEM+ and PE are other reported
comparators, but are not in this costed core. If usable DiVE code and model
weights or metric-oracle weights cannot be obtained, this strict plan is
conditional and needs a new feasibility estimate; DiME variations alone
cannot replace it.

\textbf{Replication core:} (1) DiME smile counterfactuals, 10 B200-h;
(2) young counterfactuals, 10 h; (3) one guidance, diversity, or correlation
ablation such as the paper's direct noisy-classifier guidance, 12 h; a matched
DiVE smile comparison, 8 h; shared metrics, 4 h. Total \textbf{44 B200-h $\times$ \$3.49 = \$153.56}. Alternate sampling-step or strength grids add roughly
5--12 h per setting. FID from only 100--200 counterfactuals is unstable;
report it as exploratory and foreground validity, identity, and paired
uncertainty intervals.

\subsection{DiffAtlas: Image-Mask Diffusion}
Use patient-level splits of accessible MMWHS CT/MRI heart data, five labels,
Dice and normalized surface distance. Feasible settings include same-domain
full training, two/four-shot learning, CT-to-MRI or MRI-to-CT zero-shot
transfer, per-class boundary analysis, atlas/registration versus segmenter
versus mask-diffusion baselines, and signed-distance or image-replacement
guidance ablations. A 2D or reduced-resolution model is materially different
from the paper's volumetric setting and cannot claim its original scores.

\textbf{Baselines core:} (1) nnU-Net-style segmenter, 12 B200-h;
(2) MedSegDiffv2-style mask diffusion, 24 h; (3) CMMAS registration and label
propagation on CPU, with shared scoring/inference and validation-selected
sampling steps on GPU for 4 h. Total \textbf{40 B200-h $\times$ \$3.49 = \$139.60}, \emph{plus unpriced remote CPU
time} for registration. Other feasible checks are per-class Dice/NSD and
small hyperparameter changes using those trained baselines.

\textbf{Replication core:} (1) DiffAtlas full CT training, 22 B200-h;
(2) two-shot retraining, 12 h, and CT-to-MRI zero-shot inference with Dice/NSD
and boundary plots, 6 h; (3) a targeted inference ablation of noisy-image
replacement frequency using the trained model, 12 h. Total \textbf{52 B200-h $\times$ \$3.49 = \$181.48}. Four-shot retraining or MRI-to-CT
transfer adds roughly 6--22 B200-h depending on the chosen condition. Data
download, resampling, and registration may take longer than the GPU runs.

\subsection{FLUX.1 Kontext: Flow Matching for In-Context Image Generation and Editing in Latent Space}
The released Kontext [dev] checkpoint is an image-editing model. The paper's
text-to-image and much of its headline KontextBench comparison involve
[pro]/[max] and proprietary services; [dev] alone cannot reproduce those
rankings. Feasible [dev] evaluations include
local/global/text editing, character/style references, visual-cue edits,
three- to six-turn identity drift, fixed-resolution latency, and VAE
reconstruction (PSNR/SSIM/perceptual distance). Use 50--100 fixed public
prompts per category, paired seeds, and blinded human scoring when needed.

\textbf{Supplementary editor study, not a baselines-track core:} (1) FLUX.1 Fill-dev for masked local edits,
6 B200-h; (2) HiDream-E1 for instruction edits, 6 B200-h; (3)
InstructPix2Pix on the same paired prompts, 3 B200-h; shared automated
scoring, 3 B200-h. Total \textbf{18 B200-h $\times$ \$3.49 = \$62.82}. These are accessible public editors, not an
exact replay of the paper's Figure 8 comparison. Its reported comparators
include GPT-Image-1 and Gen-4 References, which require API access and
human-preference evaluation beyond a B200 rental. A strict baselines-track
plan needs those actual comparators and a separate API/human-cost estimate.

\textbf{Replication core, evaluation-only:} (1) Kontext [dev] local and text
editing on a fixed KontextBench subset, 8 B200-h; (2) character/style subset
with paired human preference and AuraFace character-similarity scoring where
applicable, 8 h; (3) multi-turn drift with a targeted chained-previous-edit
versus original-context conditioning ablation, 6 h; VAE reconstruction and
scoring, 4 B200-h. Total \textbf{26 B200-h $\times$ \$3.49 = \$90.74}. Further feasible
checks are latency-versus-steps and reconstruction quality. If KontextBench
is not obtainable, use a public editing subset and name the benchmark change.
Paid proprietary APIs and human-review time are not in these GPU totals. This
reproduces only [dev] evaluation behavior and an inference-conditioning
ablation, not the training method or [pro]/[max] headline results.

\subsection{Qwen-Image Technical Report}
The system combines a 20B diffusion backbone with a 7B conditioner. Full
data collection, pretraining, and reinforcement learning are outside this
budget. Feasible released-weight evaluations include GenEval/DPG prompt
composition, English and Chinese short/long text accuracy using one OCR
protocol, VAE reconstruction, editing by category if the edit checkpoint is
available, and depth/novel-view tasks only if their specialized checkpoints
are available. A 12 B200-h slot could mean one B200 for twelve hours or two
for six elapsed hours if two-card capacity is available; count \emph{both}
devices. Fix seeds and prompt lists.

\textbf{Baselines core:} (1) FLUX.1-dev and SD3.5-Large on matched English
composition prompts, 16 B200-h; (2) the same two on multilingual/long-text
prompts, 16 B200-h; (3) public FLUX/Wan VAE reconstruction, 4 B200-h; shared
OCR, scoring, and a validation-selected guidance setting on the generation
baselines, 4 B200-h. Total \textbf{40 B200-h $\times$ \$3.49 = \$139.60}.
Closed comparison-model API charges are omitted, so this does not cover all
published baselines.

\textbf{Replication core, evaluation-only:} (1) released Qwen-Image
text-to-image on a fixed GenEval or DPG subset, 12 B200-h; (2) English
CVTG-2K and ChineseWord/LongText-Bench typography subsets, 12 B200-h;
(3) a targeted guidance-scale ablation on the same prompts and seeds,
12 B200-h; VAE reconstruction, 4 B200-h; shared OCR and scoring, 4 B200-h.
Total \textbf{44~\mbox{B200-h} $\times$ \$3.49 = \$153.56}. English/Chinese editing
is an optional separate evaluation only if the editing checkpoint is
available, and needs its own roughly 12 B200-h/\$41.88. A small text
benchmark cannot establish the full multilingual ranking, nor can these
released-weight tests reproduce the paper's training curriculum.

\section{Topic 3: Foundation Models, weeks 4--5}

\subsection{DINOv3}
The affordable questions concern \emph{frozen} released encoders, not
seven-billion-image pretraining or the causal effect of Gram anchoring.
Feasible evaluations include image-classification linear probes on CIFAR-100,
Pets, or an ImageNet subset; VOC/ADE20k dense segmentation; NYUv2 depth;
keypoint correspondence; TokenCut object discovery; retrieval; DAVIS mask
propagation; patch-feature PCA maps; resolution sensitivity; and matched
ViT-S/B comparisons with DINOv2. Match token count/resolution when patch
sizes differ, and use identical probe capacity and labeled splits.

\textbf{Baselines core:} (1) paper-reported DINOv2-B and SigLIP2-B frozen
classification probes, 4 B200-h; (2) matched VOC segmentation probes,
8 h; (3) matched NYUv2 depth probes, 6 h; shared extraction and validation,
including probe learning-rate/weight-decay tuning on a validation split,
2 h. Total \textbf{20 B200-h $\times$ \$3.49 = \$69.80}. Hold out the test
split until those settings are chosen. Retrieval or
patch-correspondence evaluation is another roughly 2--6 B200-h with cached
features.

\textbf{Replication core, evaluation-only:} (1) DINOv3-B versus DINOv2-B
VOC mIoU, 7 B200-h; (2) NYUv2 depth RMSE, 5 h; (3) CIFAR-100, one of the
paper's Fine-S datasets, with its frozen-feature classification protocol
and a targeted 256/512-pixel resolution ablation, 4 h; shared validation,
2 h. Total \textbf{18 B200-h $\times$ \$3.49 = \$62.82}. Keypoint or
object-discovery checks are optional. This supports transfer claims at
ViT-B scale but does not test whether Gram anchoring caused the improvement.

\subsection{V-JEPA 2.1}
The full 135k+12k self-supervised pretraining iterations and VisionMix163M
are outside the budget. Released distilled ViT-B/L features can support
reduced Something-Something-v2 or Diving48 action probes, VOC/ADE20k
segmentation, NYUv2/KITTI depth, DAVIS first-mask propagation,
last-layer-versus-four-layer probes, clip-length/resolution changes, and PCA
feature maps. Ego4D anticipation, robot grasping/navigation, and VideoQA
require substantially more task-specific data and setup than this deadline
allows.

\textbf{Baselines core:} (1) paper-reported V-JEPA 2 versus DINOv2 on a
fixed reduced Something-Something-v2 action split, 15 B200-h; (2) matched
VOC segmentation probes, 7 h; (3) DAVIS17 propagation, 7 h; feature
extraction and validation-selected probe hyperparameters, 3 h. Total
\textbf{32 B200-h $\times$ \$3.49 = \$111.68}. A frozen depth comparison adds
roughly 6--10 B200-h. The image/video tokenizer and clip protocol must be
held consistent where possible.

\textbf{Replication core, evaluation-only:} (1) Released V-JEPA 2.1
distilled-backbone Something-Something-v2 action probe with paper-matched
accuracy metric, 14 B200-h; (2) VOC dense segmentation, 6 h; (3) NYUv2
depth, 6 h; targeted last-layer-versus-four-layer or clip-length ablation,
4 h; setup, 2 h.
Total \textbf{32 B200-h $\times$ \$3.49 = \$111.68}. A reduced 256-pixel,
single-view video evaluation is modified relative to the paper's larger
resolution/multiview protocol. It does not test the dense self-supervised
training losses. Confirm that the selected checkpoint and tokenizer are
actually released before choosing this paper.

\subsection{SAM 3}
Feasible released-checkpoint tasks include text-only promptable concept
segmentation with positive and absent-class prompts on COCO/ODinW subsets;
image-only and text-plus-image exemplars; positive/negative exemplar curves;
CountBench counting; SAM 1/2 versus SAM 3 interactive point/box masks; DAVIS
video object segmentation; a small text-prompt video-concept subset; and
one-domain 0/1/3/10-shot ODinW fine-tuning. Use box/mask AP or gated F1 only
when the annotations support it. The SA-Co data engine, full ontology, and
pretraining are not feasible; substituting COCO cannot reproduce SA-Co/Gold's
reported aggregate metric.

\textbf{Baselines core:} (1) the paper-reported OWLv2 and GroundingDINO box
detectors paired with SAM 1 on a fixed released SA-Co/Gold subset if
obtainable, otherwise a disclosed LVIS phrase subset, 8 B200-h; (2) SAM 1
versus SAM 2 interactive point/box masks on matched examples, 4 h;
(3) positive/absent-class threshold tuning and hard-negative analysis of
the concept baselines, 4 h; setup, 2 h. Total \textbf{18 B200-h $\times$ \$3.49 = \$62.82}. Report concept AP or gated F1 only where annotations permit;
DAVIS VOS is an optional secondary baseline check.

\textbf{Replication core, evaluation-only:} (1) SAM 3 on matched positive
and absent-class image/phrase pairs from a released SA-Co/Gold subset if
obtainable, otherwise a named public proxy, 5 B200-h; (2) text-only,
image-only, and combined exemplars on the same positives as the targeted
prompt-modality ablation, 4 h; (3) DAVIS VOS, 8 h; CountBench
counting, 2 h; setup, 3 h. Total \textbf{22 B200-h $\times$ \$3.49 = \$76.78}. A one-domain 10-shot fine-tune can replace counting/VOS and adds
about 8--20 B200-h (\$27.92--\$69.80). Check checkpoint terms and benchmark
availability; preserve prompts, absent-class pairs, and confidence thresholds
across systems. A public proxy cannot reproduce the paper's SA-Co/Gold
aggregate score, and checkpoint inference cannot test pretraining.

\subsection{AnyUp: Universal Feature Upsampling}
This is one of the few Topic 3 papers with an affordable trainable method:
the paper reports about five hours on one H100 for 100k ImageNet steps.
Feasible checks include bilinear, guided-filter, FeatUp, and JAFAR
comparisons on VOC/COCO/ADE20k; NYUv2 depth and normals; output-resolution
sweeps; reusing a low-resolution-trained probe after upsampling; transfer
from DINOv2-S features to DINOv3/SigLIP; feature preservation; runtime and
peak memory; and attention, feature-path, self-consistency, or basis-filter
ablations. Align the frozen backbone and probe protocol across methods.

\textbf{Baselines core:} (1) the paper-reported bilinear, FeatUp, LoftUp,
and JAFAR comparators on VOC segmentation, 8 B200-h; (2) NYUv2 depth/normal
probes, 8 h; (3) resolution and
fixed-probe feature-preservation test, 6 h; setup, 2 h. Total \textbf{24 B200-h $\times$ \$3.49 = \$83.76}. Cross-backbone inference or
runtime/VRAM benchmarking can be added for roughly 2--6 B200-h.

\textbf{Replication core:} (1) One AnyUp 100k-step run, budget 8 B200-h
versus the paper's approximately five; (2) cross-backbone DINOv2-to-DINOv3
and SigLIP transfer, 6 B200-h; (3) two shortened 25k-step ablations, one
without the feature path and one without self-consistency, 3+3 B200-h;
matched bilinear/FeatUp/LoftUp/JAFAR inference, 4 B200-h; probe training/QA,
6 B200-h. Total \textbf{30 B200-h $\times$ \$3.49 = \$104.70}. More full-length ablations add about
8 B200-h/\$27.92 each. The 25k-step ablations test direction rather than
exact final performance. Verify ImageNet access and NATTEN CUDA compatibility;
if training on another image collection, declare that distribution change.

\subsection{Medical SAM Adapter (Med-SA)}
Feasible tests include BTCV organ Dice/HD95 with patient-held-out volumes;
2D REFUGE2, TNMIX, or ISIC segmentation; random one/three-point versus
approximate-overlap box prompts; SAM, U-Net/nnU-Net, or released MedSAM
comparisons; adapter rank, placement, and scale; HyP-Adpt versus additive
prompt conditioning; SD-Trans versus independent 2D slices; and prompt
perturbation. The central 3D claim requires volumetric data and SD-Trans:
a 2D-only test cannot support it. The paper uses SAM ViT-H, 1024-pixel 2D
or 128-cubed 3D inputs, 40/60 epochs on four A100s. Our smaller ViT-B plan
is explicitly a modified replication.

\textbf{Baselines core:} (1) zero-shot SAM and matched prompt sweep on a
BTCV subset, 10 B200-h; (2) the paper-reported nnU-Net volumetric baseline
at 128 cubed, 20 h; (3) released MedSAM on the same patients and prompts,
10 h; common patient-level Dice/HD95, 4 h. Total \textbf{44 B200-h $\times$ \$3.49 = \$153.56}. Match the paper's one/three-point and
0.5/0.75-overlap box conditions where supported. A 2D-only REFUGE2 fallback
could cost 40--60 B200-h (\$139.60--\$209.40) but misses the 3D conclusion.

\textbf{Replication core:} (1) Med-SA with ViT-B on 12--20 BTCV training
volumes, 24 B200-h; (2) independent-slice adapter without SD-Trans, 24 h;
(3) HyP-Adpt versus additive prompt conditioning, 24 h; shared one/three
point and two box-overlap settings, Dice/HD95, 6 h; matched zero-shot SAM
reference, 4 h; setup/QA, 6 h. Total \textbf{88 B200-h $\times$ \$3.49 = \$307.12}. Adapter-placement or rank
adds one training condition, roughly 24 B200-h/\$83.76. Keep volume-level splits
and generate prompts from the correct masks; reductions in backbone size,
patients, and steps limit comparison with original 17-task claims.

\subsection{Towards a Unified View of Parameter-Efficient Transfer Learning}
Feasible design-axis tests use SST-2/MNLI classification, XSum summarization,
or WMT16 English-Romanian translation. Compare full tuning, BitFit, prompt,
prefix, sequential/parallel adapters, LoRA, and Mix-And-Match (MAM) at matched
\emph{trainable parameter counts}; examine attention versus FFN insertion,
parallel-adapter scale, LoRA scale, prefix length, and transfer to a second
task. Report accuracy/ROUGE/BLEU and tuned parameter fraction. SST-2 alone
is too easy to support the paper's broader architecture ranking.

\textbf{Baselines core:} (1) RoBERTa-base SST-2 full/BitFit/prefix/adapter/
LoRA comparison, 8 B200-h; (2) BART-large on a fixed 20k-example XSum
subset with prefix, sequential/parallel adapter, and LoRA, 30 h; (3)
matched-parameter rank/bottleneck sweep, 10 h; shared scoring, 4 h. Total \textbf{52 B200-h $\times$ \$3.49 = \$181.48}. Reduced MNLI or an extra
decoding-setting check is feasible within roughly 10--20 more B200-h.

\textbf{Replication core:} Seven approximately six-hour BART-large XSum
runs cover three design axes with an explicit matched control: (1) full tuning
and sequential versus parallel FFN adapters; (2) parallel attention versus
FFN insertion; (3) LoRA-FFN scale one versus four and MAM (prefix length 30,
parallel adapter rank 512). The seven conditions are full tuning,
sequential-FFN adapter, parallel-FFN adapter, parallel-attention adapter,
LoRA scale one, LoRA scale four, and MAM.
Scoring/QA adds 6 B200-h. Total \textbf{48 B200-h $\times$ \$3.49 = \$167.52}. A second 10k-example English-Romanian check
adds roughly 20--50 B200-h (\$69.80--\$174.50). The 20k XSum subset, fewer seeds,
and shorter training differ from the paper's much larger setup; equalize
trainable parameters rather than merely setting equal rank values.

\end{document}
